\scanpage{147}
\mspage{21}
je vais d'abord majorer le cardinal des classes d'isomorphie d'objets de $M^r_{/X}$, i.e. des \deleted{objets} flèches à isom. près (dans $M^r$)
\diagram{\begin{tikzcd}[ampersand replacement=\&,column sep=large]
A\arrow[r]\&X
\end{tikzcd}}
\margin{$X$ fixé dans $\mathrm{Filt}^s$\quad $A$ variable dans $M^r$.}
On aura
\[
A=\varinjlim_{i\in I}A_i\qquad
\begin{aligned}
&\text{avec }A_i\in\mathrm{Filt}^{\pi_0}\\
&I\text{ ens ordonné de cardinal }\le r^{\pi_0}
\end{aligned}
\]
Pour $I$ fixé, le \deleted{nombre des} cardinal de l'ensemble des systèmes inductifs indexés par $I$, dans $M^{\pi_0}_{/X}$, \deleted{est majoré par} à isom près, peut se majorer ainsi : pour \deleted{tt} $i\in I$, il faut choisir $A_i$ \inserted{dans $M^{\pi_0}$} (\deleted{dans} \unsure{seul|seule} sa classe dans $\widetilde{M^{\pi_0}}$ importe), donc choix majoré par \deleted{$\unsure{C_0|C_1}$} $\mathrm{card}\,\widetilde{M^{\pi_0}}\inserted{\le C_0}$, et pour tous les $i$, par
\[
\mathrm{card}\,I\times C_0\le r^{\pi_0}\times C_0,
\]
et \deleted{comme} \inserted{si} $r\ge C_0$, \inserted{donc $r^{\pi_0}\ge C_0$,} on a donc une majoration \deleted{\unsure{des|du}} choix possibles d'un système des \uline{objets} $A_i$ \deleted{$\unsure{i\in I|\pi}$} dans $M^\pi$, à isom près, par $r^{\pi_0}$.
% slanted
\margin{On a \unsure{pris ici} $C_0=\mathrm{card}\,\widetilde{\underline{\mathrm{Fl}}^{\pi_0}}$, donc $C_0\le C_1$}
Si on se donne de plus les morphismes de transition $A_i\mathbin{-}A_j$ pour $i\unsure{\subset|<}j$, cela fait pour chaque paire \unsure{compatibles|comparables}, \inserted{\unsure{élém.|élémts|élus}} dans $I\times I$, un choix majoré par $C_0$, donc encore une majoration par
\[
\underbrace{\mathrm{card}(I\times I)}_{=(\mathrm{card}\,I)^2=\mathrm{card}\,I}\times C_0\le r^{\pi_0}\cdot C_0=r^{\pi_0}.
\]
Donc \inserted{l'ens.} \unsure{des|les} systèmes inductifs en question, à isom près, pour $I$ fixé, sont de cardinal
\[
\le r^{\pi_0}\cdot r^{\pi_0}=r^{\pi_0}.
\]
Enfin, les \uline{\unsure{hom.|lim.}} d'un syst. inductif donné dans $X$ sont majorés par
\[
\mathrm{card}\,I\times\unsure{s|6}^{(r^{\pi_0})}
\]
(par 3.8. déjà prouvé), donc par
\[
r^{\pi_0}\cdot s^{(r^{\pi_0})}=s^{(r^{\pi_0})}\qquad(\text{puisque }r^{\pi_0}\unsure{\le|<}s^{(r^{\pi_0})}).
\]
\deleted{\unsure{Donc|D'onc} \unsure{Au total, pour $I$|pour $I$}} \inserted{Au total, pour $I$ fixé,} \inserted{l'ens.} des syst. inductifs à isom près, indexés par $I$, \unsure{à val.} dans $M^r_{/X}$, \deleted{\unsure{sont}} \unsure{\deleted{est}|est}
% archive sheet number: 146
